\section{Examples}
\label{sec:examples}

We describe two functional records from the corpus. Both concern well-known,
publicly disclosed and long-patched vulnerabilities; the point of the records
is not the vulnerabilities themselves, but the reusable, verified datasets
they produce. The first, a local privilege escalation in \texttt{sudo}, was
migrated by hand from the previous generation of \cvex and is verified by a
planted flag. The second, an HTTP request-smuggling flaw in the
\texttt{aiohttp} Python framework, was converted from a \genie reproduction
and is verified by \genie's CTF-style verifier. Table~\ref{tab:examples}
summarizes them.

\begin{table}[t]
\centering
\caption{The two example records.}
\label{tab:examples}
\small
\begin{tabular}{@{}lll@{}}
\toprule
 & \textbf{CVE-2023-22809} & \textbf{CVE-2024-52304} \\
\midrule
Software & \texttt{sudo} 1.9.12p1 & \texttt{aiohttp} 3.10.10 \\
CWE & CWE-269 & CWE-444 \\
Class & local privilege escalation & HTTP request smuggling \\
Origin & previous \cvex generation & \genie \\
Topology & \texttt{single} & \texttt{single} \\
Success criterion & planted flag & \genie verifier \\
Flag location & \texttt{/root/flag} (mode 0600) & --- \\
Traced process & \texttt{sudo} & \texttt{python} \\
Attack capture & 2,164 syscalls & 515~KB syscalls, 4.5~KB pcap \\
\bottomrule
\end{tabular}
\end{table}

\subsection{Local: sudoedit (CVE-2023-22809)}

\paragraph{The vulnerability.}
\texttt{sudoedit} (\texttt{sudo -e}) lets an administrator allow a user to
edit specific files with elevated privileges, without granting a root shell.
According to the public advisory, versions of \texttt{sudo} up to 1.9.12p1
mishandle the editor selected through the user's environment, so that a user
who is permitted to edit one file can edit other files as
well~\cite{cve}. Since the user controls which files are opened with
elevated privileges, the flaw amounts to a full privilege escalation. It was
fixed in \texttt{sudo} 1.9.12p2.

\paragraph{The target.}
The target image builds the vulnerable \texttt{sudo} 1.9.12p1 from a
vendored source tarball, so that the record does not depend on a
distribution package that may change or disappear. It creates an
unprivileged user and installs a sudoers policy that grants this user a
single right: to edit \file{/etc/motd} with \texttt{sudoedit}. At container
start, \file{plant-flag.sh} writes the record's flag to \file{/root/flag},
owned by root with mode 0600. The policy is the whole point of the record:
under a correct \texttt{sudo}, nothing the user is allowed to do can reveal
the flag.

\paragraph{Benign and attack phases.}
The benign phase exercises exactly the permission the policy grants: the
user edits \file{/etc/motd} through \texttt{sudoedit}, with an editor that
leaves the file unchanged. The attack phase starts from a fresh container
and runs, as the same unprivileged user, the publicly documented
exploitation of the flaw; it succeeds if it prints the flag. Both phases run
the same binary under the same policy, so the two traces differ only in what
the exploitation adds.

\paragraph{Verification and dataset.}
The record declares \texttt{success: flag} with the \texttt{stdout} channel:
the runner compares the flag printed by the attack with the planted one, and
in the reference run they match. Because the flag is readable only by root,
a match proves that the escalation genuinely happened, and not merely that
an exploit ran. The record is local, so its dataset is centered on system
calls: the recorder traces \texttt{sudo} and its children, producing 1,699
lines in the benign phase and 2,164 in the attack phase. The difference is
visible directly in the traces: the attack trace contains the access to
\file{/root/flag} and additional process executions, while the benign trace
never touches the flag. This pairing (same program, same policy, one benign
and one malicious use) is the kind of labeled data that host-based detectors
need.

\subsection{Network-Facing: aiohttp (CVE-2024-52304)}

\paragraph{The vulnerability.}
\texttt{aiohttp} is a widely used asynchronous HTTP client/server framework
for Python. According to its advisory, before version 3.10.11 the
pure-Python HTTP parser incorrectly handled newlines inside chunk extensions
of chunked request bodies~\cite{cve}. When the C extensions are not in use,
a front-end proxy and an \texttt{aiohttp} back end can disagree on where a
request ends, which enables HTTP request smuggling (CWE-444) and can let an
attacker bypass protections enforced by the proxy.

\paragraph{From a \genie reproduction to a record.}
\genie reproduced this CVE automatically, producing an exploit and a
verifier. The verifier forces the pure-Python parser, starts an
\texttt{aiohttp} server, runs the exploit against it, and checks whether the
server received two distinct requests although the client sent only one;
if so, it prints \genie's flag. The conversion to \cvexthree packages these
artifacts into a record: the target image is \texttt{python:3.11-slim} with
\texttt{aiohttp} 3.10.10 (the last release before the fix) installed from
PyPI, and the exploit and verifier are copied into it. The record's
\file{\_source} directory keeps \genie's original artifacts and the CVE
metadata (advisory, CWE, patch commit), so that the record's provenance is
preserved.

\paragraph{Verification and dataset.}
The record declares \texttt{success: verifier}: the attack phase runs
\genie's verifier inside the target, and the runner accepts the
reproduction if the verifier succeeds; in the reference run it prints the
flag \texttt{3xploit66full}. Because the verifier starts its own server and
talks to it over the loopback interface, the recorder, which shares the
target's network namespace, captures the smuggling exchange itself (a
4.5~KB packet capture) together with 515~KB of system-call trace of the
Python processes. The record is a network-facing vulnerability reproduced
within a single container: no separate attacker container is needed, yet
the dataset still contains the malicious HTTP traffic.

\paragraph{Limitations.}
This record also shows where the corpus still needs work. Its benign phase
only imports the library, a placeholder rather than a realistic workload;
for most records converted from \genie, the benign workloads have yet to be
written so that they exercise the same interface as the attack (here, normal
chunked requests to the same server). Records converted from \genie also
download their source at build time rather than vendoring it. Both are on
our list of improvements to the corpus.
